\section{VAE-IW}

Although width-based planning has been shown to be generally very effective for classical planning domains \cite{iworginal,BFSW,ijcai2017-600}, its performance in practice significantly depends on the quality of the features used. Features that better capture meaningful structure in the environment typically translate to better planning results \cite{RLIW}. 
Here we propose VAE-IW, in which representations extracted by a VAE are used as features for RolloutIW. 
The main advantages are:
\begin{itemize}
    \item Significantly fewer features than B-PROST, leading to faster planning and a smaller memory footprint.
    \item Autoencoders, in particular VAEs, are a natural approach for learning meaningful, compact representations from data \cite{bengio2013representation,tschannen2018recent,kingma2019introduction}.
    \item No additional preprocessing is needed, such as background masking~\cite{rolloutiw-orginal}.
\end{itemize}
On the planning side, we follow \citet{RLIW} and use RolloutIW($k$) with width $k=1$ to keep planning efficient even with a small time budget.
% This is sufficient to obtain good results, and at the same time makes it easier to compare our experimental results with the aforementioned paper. 
In the remainder of this section, we detail the two main components of the proposed method: feature learning and planning.

\paragraph{Unsupervised feature learning.}

For feature learning we use a variational autoencoder with a discrete latent space:
\begin{equation}
    \ptheta(\xb) = \sum\nolimits_{\zb} \ptheta(\xb, \zb) = \sum\nolimits_{\zb} \ptheta(\xb \given \zb) p(\zb)
\label{eq:vaeiw-marginal}
\end{equation}
where the prior is a product of independent Bernoulli distributions with a fixed parameter:
% \begin{equation}
%     p(\zb) = \prod_{i=1}^K p(z_i) = \prod_{i=1}^K \Bernoulli(\mu) \ .
% \end{equation}
    $p(\zb) = \prod_{i=1}^K p(z_i)$, where $ p(z_i) = \Bernoulli(\mu)$.
Given an image $\xb$, we would like to retrieve the binary latent factors $\zb$ that generated it, and use them as a propositional representation for planning. Since $\ptheta(\xb \given \zb)$ is parameterized by a neural network, and therefore highly nonlinear with respect to the latent variables, inference of the latent variables is intractable (it would require the computation of the sum in~\eqref{eq:vaeiw-marginal}, which has a number of terms exponential in the number of latent variables).

We define an inference model:
% \begin{align}
%     \qphi(\zb \given \xb) 
%     &= \prod_{i=1}^K \qphi(z_i \given \xb) \\
%     &= \prod_{i=1}^K \Bernoulli((\mu_\phi(\xb))_i)
% \end{align}
\begin{align*}
    \qphi(\zb \given \xb) 
    &= \prod_{i=1}^K \qphi(z_i \given \xb) 
    = \prod_{i=1}^K \Bernoulli((\mu_\phi(\xb))_i)
\end{align*}
to approximate the true posterior $\ptheta(\zb \given \xb)$.
The \emph{encoder} $\mu_\phi$ is a deep neural network with parameters $\phi$ that outputs the approximate posterior probabilities of all latent variables given an image $\xb$.
Using stochastic amortized variational inference, we train the inference and generative models end-to-end by maximizing the ELBO with respect to the parameters of both models.
We approximate the discrete Bernoulli distribution of $\zb$ with the Gumbel-Softmax relaxation~\cite{jang2016categorical,maddison2016concrete}, which is reparameterizable. This allows us to estimate the gradients of the inference network with the pathwise gradient estimator.
Alternatively, the approximate posterior could be optimized directly using a score-function estimator~\cite{williams1992simple} with appropriate variance reduction techniques or alternative variational objectives~\cite{Bornschein2014-my,Mnih2016-vr,Le2018-kj,Masrani2019-nq,lievin2020optimal}.

Note that we are not interested in the best generative model---e.g. in terms of marginal log likelihood of the data or visual quality of the generated samples---as much as in representations that are useful for planning.
Thus, ideally, we would directly optimize $\beta$ for planning performance. In practice, however, we assume that a necessary (but not sufficient) condition for a representation to be useful in our case is that it contains enough information about the entities on the screen.
Following \citet{burgess2018understanding}, we control the tradeoff between the competing terms in~\eqref{eq:elbo} by using~\eqref{eq:beta_objective_function} and decreasing $\beta$ until good quality reconstructions are obtained.


\paragraph{Planning with learned features.}

After training, the inference model $\qphi(\zb \given \xb)$ yields approximate posterior distributions of the latent variables. The binary features for downstream planning can be obtained by sampling from these distributions or by deterministically thresholding their probabilities given a fixed threshold $\lambda \in (0, 1)$.
% \begin{equation*}
%     f_i = 
%         \begin{cases}
%           1 &\text{if}\ \mu_\phi(\xb)_i \geq \lambda \\
%           0 &\text{otherwise}\\
%         \end{cases}
% \end{equation*}
% where $\lambda \in (0, 1)$ is the threshold.
% In this work we choose to deterministically threshold the probabilities, as it empirically yields more stable features for planning.
In this work we choose the latter approach as it empirically yields more stable features for planning.
Overall, this provides a way of efficiently computing a compact, binary representation of an image, which can in turn be interpreted as a set of propositional features to be used in symbolic reasoning systems like IW or RolloutIW. 

The planning phase in VAE-IW is based on RolloutIW, but uses the features extracted by the inference network instead of the B-PROST features. Note that while the B-PROST features include temporal information, the VAE features are computed independently for each frame. This should, everything else being equal, give an advantage to planners that rely on B-PROST features.
